% !TEX root = ../notes_missingsemester.tex
% Block 3: Package the Stack — Dependencies & Project Structure
%%%%%%%%%%%%%%%%%%%%%%%%%%%%%%%%%%%%%%%%%%%%%%%%%%%%%%%%%%%%%%%%%%%%%%

\index{dependencies}\index{virtualenv}\index{pip}\index{project structure}
\chapter{Package the Stack: Dependencies \& Project Structure}
\newthought{Reproducibility starts before you write a single line of application code.}

% ==============================================================================
% SESSION A: THEORY
% ==============================================================================

\section{The Why}

\newthought{``Works on my machine''} is the most dangerous sentence in software engineering.
Your colleague clones the repo, runs the code, and gets an \texttt{ImportError}. You installed
\texttt{scikit-learn} months ago and forgot. They did not. The code is correct; the environment is wrong.

Dependency management solves this: declare what your project needs, pin the exact versions, and
anyone can reconstruct the identical environment in two commands.

\section{Virtual Environments}\index{virtualenv}\index{venv}

\newthought{A virtual environment} is an isolated Python installation. Each project gets its own
set of packages, independent of every other project and the system Python.

\begin{verbatim}
python3 -m venv .venv
source .venv/bin/activate
\end{verbatim}

\marginnote{The \texttt{.venv/} directory contains a full Python installation: interpreter, \texttt{pip},
and all installed packages. It is machine-specific and large (often 50–200\,MB).
It must be in \texttt{.gitignore}: never commit it. Anyone can recreate it from \texttt{requirements.txt}.}

\newthought{While the environment is active,} \texttt{python} and \texttt{pip} point to the local copies.
Install a package and it stays inside \texttt{.venv/}, with no system-wide pollution. Deactivate with \texttt{deactivate}.

\subsection{One Project, One Environment}

\newthought{The rule is simple:} every project gets its own \texttt{.venv}. Project A needs
\texttt{numpy==1.24}; Project B needs \texttt{numpy==2.0}. Without isolation, one of them breaks.
With virtual environments, they coexist peacefully on the same machine.

\section{Managing Packages with pip}\index{pip}\index{requirements.txt}

\newthought{pip is the package installer} for Python. The core cycle:

\begin{verbatim}
pip install scikit-learn joblib python-dotenv
pip freeze > requirements.txt
pip install -r requirements.txt
\end{verbatim}

\begin{description}
    \item[\texttt{pip install}] Downloads and installs packages from PyPI\index{PyPI} (the Python Package Index):
    \url{https://pypi.org/}.
    \item[\texttt{pip freeze}] Lists every installed package with its exact version.
    \item[\texttt{pip install -r}] Installs everything listed in a requirements file—the reproducibility step.
\end{description}

\subsection{Anatomy of \texttt{requirements.txt}}

\newthought{One package per line,} pinned to an exact version:

\begin{verbatim}
scikit-learn==1.4.0
joblib==1.3.2
python-dotenv==1.0.1
\end{verbatim}

\marginnote{Version specifiers: \texttt{==} pins exactly, \texttt{>=} sets a floor, \texttt{\~{}=} allows
patch updates (\texttt{\~{}=1.4.0} means \texttt{>=1.4.0, <1.5.0}). For maximum reproducibility, pin with \texttt{==}.}

\section{Semantic Versioning}\index{semantic versioning}

\newthought{Version numbers are not arbitrary.} The convention \texttt{MAJOR.MINOR.PATCH} signals
what changed:

\begin{description}
    \item[MAJOR] Breaking changes. Your code may need to be rewritten. Upgrading is a deliberate decision.
    \item[MINOR] New features, backwards-compatible. Safe to upgrade, but test.
    \item[PATCH] Bug fixes only. Safe to upgrade.
\end{description}

\newthought{When \texttt{scikit-learn} goes from 1.4.0 to 2.0.0,} pay attention. When it goes from 1.4.0
to 1.4.1, relax. This convention lets you make informed upgrade decisions without reading every changelog.

\section{Configuration: \texttt{.env} and \texttt{python-dotenv}}\index{.env}\index{python-dotenv}\index{12-factor app}

\newthought{Separating configuration from code} is a foundational principle. The 12-factor app methodology\marginnote{The
Twelve-Factor App: \url{https://12factor.net/}\\
python-dotenv: \url{https://github.com/theskumar/python-dotenv}}
states: store config in the environment, never in the source code.

\newthought{A \texttt{.env} file} holds key-value pairs that your application reads at startup:

\begin{verbatim}
SECRET_API_KEY=my-actual-secret-key-here
DEBUG=true
\end{verbatim}

In Python, \texttt{python-dotenv}\index{python-dotenv} loads these into environment variables:

\begin{verbatim}
from dotenv import load_dotenv
import os

load_dotenv()
api_key = os.getenv("SECRET_API_KEY")
debug = os.getenv("DEBUG", "false").lower() == "true"
\end{verbatim}

\subsection{The \texttt{.env.example} Pattern}\index{.env.example}

\newthought{The \texttt{.env} file contains real secrets}—it must never be committed.
The \texttt{.env.example} file is the contract: it shows which variables exist, what they are for,
and provides safe placeholder values:

\begin{verbatim}
# API key callers must send in the X-API-Key header
# (a secret — never a real value here)
SECRET_API_KEY=replace-me

# Debug mode: when true, FastAPI shows /docs and full
# error tracebacks. Never true in production.
DEBUG=true
\end{verbatim}

\marginnote{\texttt{.env.example} is committed to Git and public. \texttt{.env} is in \texttt{.gitignore}
and stays local. The contrast is immediate: one is the template, the other is the secret.
This pattern grows block by block—every new service adds its variables to \texttt{.env.example}.}

Students create their actual \texttt{.env} by copying \texttt{.env.example} and replacing the placeholders
with real values.

\section{Security Auditing with \texttt{pip audit}}\index{pip audit}\index{CVE}

\newthought{Your dependencies have dependencies,} and any of them can have known security vulnerabilities.
\texttt{pip audit} checks your installed packages against a database of known CVEs\index{CVE}
(Common Vulnerabilities and Exposures):

\begin{verbatim}
pip install pip-audit
pip-audit
\end{verbatim}

\marginnote{A \textbf{CVE} is a publicly disclosed security vulnerability with a unique identifier
(e.g., CVE-2024-12345). The National Vulnerability Database (\url{https://nvd.nist.gov/}) tracks them.
\texttt{pip-audit}: \url{https://github.com/pypa/pip-audit}}

If vulnerabilities are found, the output tells you which package, which version, and what the fix is
(usually: upgrade to a newer version).

\section{The Container Alternative}\index{Docker}

\marginnote{Docker\index{Docker} and Docker Compose solve the ``reproducible environment'' problem at the
OS level rather than the Python level. This course teaches the virtualenv path to keep the OS layer visible
and learnable. Docker is the natural next step after this course.\\
Docker: \url{https://www.docker.com/}\\
Docker Docs: \url{https://docs.docker.com/get-started/}}

\newthought{Virtual environments isolate Python packages.} Docker containers isolate the entire operating system:
libraries, system packages, Python version, everything. Where \texttt{requirements.txt} says ``install these
Python packages,'' a \texttt{Dockerfile} says ``start from this Linux image, install these system packages,
then install these Python packages.'' The result is a self-contained unit that runs identically on any machine
with Docker installed. We do not use Docker in this course—but you should know it exists and why teams adopt it.

\section{Optional: When \texttt{requirements.txt} Starts to Hurt}\index{pyproject.toml}

\marginnote{Tools like \texttt{uv}\index{uv} and \texttt{poetry}\index{poetry} use this model and are fast
becoming the standard. Once the \texttt{requirements.txt} workflow is second nature, this is the upgrade.
See \url{https://packaging.python.org/en/latest/}.}

\newthought{\texttt{pip freeze > requirements.txt}} captures everything—your direct dependencies, their
dependencies, and their dependencies' dependencies, all flattened into one list. Six months later you want
to upgrade one package: which of those 47 lines do you actually own? \texttt{pyproject.toml} solves this
by separating \emph{what you depend on} (your direct deps, loosely pinned) from \emph{what gets installed}
(a generated lock file with the full resolved graph). Upgrades become surgical: bump one line, regenerate
the lock, done.

% ==============================================================================
% SESSION B: HANDS-ON
% ==============================================================================

\newpage
\section{Workshop}

\marginnote{If your VM is broken, restore the \texttt{B2-complete} snapshot.
At the end of this session, take a snapshot named \texttt{B3-complete}.}

\begin{enumerate}
    \item \textbf{Create and activate a virtual environment} on the \texttt{dev} VM:
    \begin{verbatim}
    cd /home/student/pixelwise
    python3 -m venv .venv
    source .venv/bin/activate
    \end{verbatim}
    Verify: \texttt{which python} should point to \texttt{.venv/bin/python}.

    \item \textbf{Scaffold the project directories.} Create only what we need right now—other folders
    appear in the block that first needs them:
    \begin{verbatim}
    mkdir -p app data models
    touch app/__init__.py
    \end{verbatim}

    \item \textbf{Install dependencies and pin them:}
    \begin{verbatim}
    pip install scikit-learn joblib python-dotenv
    pip freeze > requirements.txt
    \end{verbatim}
    Open \texttt{requirements.txt} and inspect: note the pinned versions and the transitive dependencies
    that were pulled in automatically.

    \item \textbf{Write \texttt{.env.example}:} Create the file with two entries:
    \begin{verbatim}
    # API key (secret — never a real value here)
    SECRET_API_KEY=replace-me

    # Debug mode (never true in production)
    DEBUG=true
    \end{verbatim}
    Then create your actual \texttt{.env} with a real value for \texttt{SECRET\_API\_KEY}.

    \item \textbf{Update \texttt{.gitignore}:} Add the entries that now matter—without them,
    \texttt{git status} is full of noise and secrets are one careless \texttt{git add .} away from GitHub:
    \begin{verbatim}
    .venv/
    *.pkl
    data/
    .env
    \end{verbatim}
    Run \texttt{git status} before and after—see the difference.

    \item \textbf{Run \texttt{pip audit}:}
    \begin{verbatim}
    pip install pip-audit
    pip-audit
    \end{verbatim}
    Inspect and discuss any findings. What is a CVE? What would you do if a critical vulnerability
    is found in one of your dependencies?

    \item \textbf{Commit the scaffold:}
    \begin{verbatim}
    git add app/ requirements.txt .env.example .gitignore
    git commit -m "Add project scaffold, deps, and env template"
    git push
    \end{verbatim}

    \item \textbf{Replicate on the server:} SSH into the server, pull the latest code, and recreate
    the environment:
    \begin{verbatim}
    cd /opt/pixelwise
    git pull
    python3 -m venv .venv
    source .venv/bin/activate
    pip install -r requirements.txt
    \end{verbatim}
    The two-command reproducibility promise: \texttt{venv} + \texttt{pip install -r}.

    \item \textbf{Take a snapshot} of both VMs. Name them \texttt{B3-complete}.
\end{enumerate}

\section{Security Thread}

\newthought{\texttt{pip audit}\index{pip audit}} catches known CVEs in dependencies—run it regularly and before
every deploy. \texttt{.env.example} makes the secrets contract explicit without exposing values.
Never commit \texttt{.env}. The \texttt{.venv/} folder must be in \texttt{.gitignore}—it is large,
machine-specific, and reconstructable from \texttt{requirements.txt} alone.

\section{Self-Reflection and Recap}

\newthought{Reflection questions:}
\begin{itemize}
    \item Why do we use a virtual environment instead of installing packages system-wide?
    \item What is the difference between \texttt{.env} and \texttt{.env.example}? Why do we need both?
    \item What does \texttt{pip freeze} capture that you did not explicitly install?
    \item If a colleague clones your repo, what two commands do they run to get a working environment?
    \item What is the risk of committing \texttt{.env} to Git—even briefly?
\end{itemize}

\newthought{Milestone:} PixelWise has \texttt{app/}, \texttt{data/}, \texttt{models/}, a virtualenv,
pinned dependencies, and a minimal \texttt{.env.example}. The \texttt{.gitignore} is now doing real work.
Anyone who clones the repo can recreate the environment in two commands.
